\subsection{Uncapped electronic-barrier differences and archived physical records}
\label{sec:si-archived-barrier-errors}

The positive tiers contain 36 tasks with 53 numerical reference edges per workflow. HA1, HA6, HB1, and HF1 score qualitative channel preference; NB1\_E2 and NQ5\_E2 have no numerical reference. The Impossible tier is kept separate, including IB9's one numerical edge and its negative target. We recovered 160 explicit log-derived-versus-Gold comparison notes from the original benchmark databases, comprising 159 positive-target observations and one Guided/IB9 observation. Native TS and reference energies, their echoed methods, and file hashes are retained with the accompanying analysis.

For the positive targets, 156 observations were reproduced from normally terminated native logs using the selected TS and reference files, consistent methods and continuum settings, and agreement with the original comparison note at its recorded 0.001-kcal~mol$^{-1}$ precision. Three further observations belong to runs whose original score components differ from the final publication components; excluding these leaves 153 final-aligned observations: Guided 41, Standard 37, Record 40, and Brief 35, from 33, 29, 32, and 30 tasks, respectively. The 59 remaining opportunities out of 212 positive numerical run--edge pairs remain listed as unavailable or unresolved observations. Missing values receive neither zero error nor an imputed numerical penalty. The separate final F4 replay records are outside this uniformly sourced original-note cohort.

For an observed pair, we calculate the signed difference as $(E_{\mathrm{TS}}-E_{\mathrm{ref}})\,627.509474-\Delta E^{\ddagger}_{\mathrm{Gold}}$, in kcal~mol$^{-1}$, using native Hartree energies and the frozen task-file reference. Mean absolute error (MAE), root-mean-square error (RMSE), median absolute error, and maximum absolute error have no Gaussian transformation or penalty cap. Each observed reference edge has equal weight.

The intersection shared by all four workflows contains 26 reference edges from 23 positive tasks, giving 104 run--edge observations (Table~\ref{tab:si-uncapped-common}). Their native TS and reference pairs use B3LYP-D3(BJ)/def2-SVP with no explicit continuum keyword in the original input routes. These are conditional numerical comparisons on the same recoverable edges. They retain the archived matcher's target assignments, including the local-reference L2 edge, and describe neither unobserved numerical targets nor complete reaction-path support. Guided/G2 is excluded because its final structured delivery omits the numerical barrier and TS energy/log bindings, although its historical native calculations support an approximately 0.845-kcal~mol$^{-1}$ value.

\begin{table}[!htbp]
\centering
\small
\caption{Uncapped electronic-barrier differences on the 26 reference edges shared by all four workflows, covering 23 positive tasks. All entries are in kcal~mol$^{-1}$; every workflow uses the same edge set.}
\label{tab:si-uncapped-common}
\begin{tabular}{lrrrr}
\toprule
Workflow & MAE & RMSE & Median absolute error & Maximum absolute error \\
\midrule
Guided & 1.1676 & 2.8566 & 0.0573 & 10.7826 \\
Standard & 1.8995 & 5.0883 & 0.0574 & 21.4713 \\
Record & 1.5763 & 5.7082 & 0.0731 & 28.6373 \\
Brief & 1.1310 & 2.3080 & 0.0985 & 9.0486 \\
\bottomrule
\end{tabular}
\end{table}

Retained selected Gold input files corroborate the same method and continuum setting for a further-defined common subset of 19 edges from 17 tasks (Table~\ref{tab:si-uncapped-gold-input}). This corroboration concerns the selected inputs; complete reference-log and endpoint-chain coverage remains a separate question. Record has the lowest MAE in this subset. The changed ordering shows that the numerical ranking depends on which archived reference comparisons are available.

\begin{table}[!htbp]
\centering
\small
\caption{Uncapped differences for the 19 common reference edges with additional corroboration from retained selected Gold input files, covering 17 positive tasks. Units and statistics follow Table~\ref{tab:si-uncapped-common}.}
\label{tab:si-uncapped-gold-input}
\begin{tabular}{lrrrr}
\toprule
Workflow & MAE & RMSE & Median absolute error & Maximum absolute error \\
\midrule
Guided & 1.0105 & 2.7118 & 0.0388 & 10.7826 \\
Standard & 1.2903 & 3.3038 & 0.0424 & 10.7826 \\
Record & 0.3934 & 0.8609 & 0.0424 & 3.0918 \\
Brief & 1.2811 & 2.6213 & 0.0744 & 9.0486 \\
\bottomrule
\end{tabular}
\end{table}

The archive also permits an explicit inventory of physical-record availability across the 160 positive runs. Guided, Record, and Brief each retain delivery-DAG element states for all 40 tasks; Standard's configuration has no delivery-DAG objects. Frozen delivery-disposition reports exist for Guided 28, Record 39, Brief 5, and Standard 0 runs, totalling 72/160. Reports containing native scientific checks with a graph hash matching the archived final root graph exist for Guided 28 and Brief 5 runs, totalling 33/160. These records concern the delivery scope selected during a run. Historical binary self-consistency checks cover delivered graph elements, and a DAG completion report can include documented skipped or not-applicable stages. A common rate of fully supported requested mechanisms therefore requires a uniform offline review of requested-object coverage and the existing native physical conditions in the final organized graphs, keeping additional exploratory branches separate. The present analysis reports the available records and numerical observations; the frozen component and composite scores remain unchanged.

The companion directory \path{companion_data/revision_20261101/metrics/} provides every numerical opportunity, original comparison note, recovered native energy and method, source/version relation, exclusion, and cohort membership. Its reproduction script recomputes the statistics from the stored native energies and frozen Gold values.
